\section*{Response to Reviewer Comments --- FHOPS
(SOFTX-D-26-00697R2)}\label{response-to-reviewer-comments-fhops-softx-d-26-00697r2}

\textbf{Manuscript:} FHOPS: a reproducible optimisation and benchmarking
stack for forest harvest operations \textbf{Decision:} Minor revision
\textbf{Editor:} Randall Sobie (Editor-in-Chief, SoftwareX)

\begin{center}\rule{0.5\linewidth}{0.5pt}\end{center}

\subsection*{Summary of changes}\label{summary-of-changes}

We thank the editor and both reviewers for their careful reading and
positive assessment. In this revision we have:

\begin{enumerate}
\def\labelenumi{\arabic{enumi}.}
\tightlist
\item
  Clarified the planning scope of the operational model in Section 2:
  where it sits in the planning hierarchy, and that haul transport is
  outside its boundary (R2.5, R2.9).
\item
  Explained how production rates, availability, and stochastic playback
  represent utilisation and operational, mechanical, and weather delays,
  and added an explicit limitation on delay costing (R2.7).
\item
  Described how head-start buffers support comparing coupled and
  decoupled harvesting systems (R2.6), and how the rolling-horizon
  driver supports re-solving after disruptions (R2.8).
\item
  Revised the keywords (R2.2), spelled out FHOPS in the abstract (R2.3),
  and improved the legibility of the playback figure (R2.10).
\item
  Moved the operational MILP formulation out of the main text, as the
  editor suggested (E.1): the full labelled formulation is versioned
  with FHOPS 1.0.1 and cited, and a compact appendix summarises it with
  an equation-to-code table.
\item
  Corrected defects in FHOPS that the review led us to find, released
  FHOPS 1.0.1, and regenerated every benchmark, tuning, playback, and
  scaling result on that release (see the note on FHOPS 1.0.1 at the end
  of this letter).
\end{enumerate}

Reviewer \#2's annotations were made on the PDF of the original
submission. Below we refer to sections of the revised manuscript.

\begin{center}\rule{0.5\linewidth}{0.5pt}\end{center}

\subsection*{Editor}\label{editor}

\subsubsection{E.1 ``You might consider putting some of the equations in
an
appendix.''}\label{e.1-you-might-consider-putting-some-of-the-equations-in-an-appendix.}

\textbf{Response:} Done. The complete formulation of the operational
MILP as implemented in FHOPS 1.0.1 (sets, parameters, decision
variables, every constraint equation with its exact linearisation, and
the list of changes from FHOPS 1.0.0) is now maintained and versioned
with the software: it is part of the FHOPS repository at tag v1.0.1 and
is rendered in the FHOPS documentation (``Operational MILP
formulation''); both are cited in Section 2.3 and Appendix A. Appendix A
is a compact summary of about four pages: brief notation, the objective
(OBJ), one sentence on the optional earliness tie-break (OBJ2), and a
table that gives, for each labelled block (E1--E13 with E8a--E8c, the
optional initial state INIT, and the domains D1), its meaning and the
Pyomo components that implement it. The labels are those of the
versioned formulation, so every equation can be traced from the paper to
the full statement and to the code. Blocks that persist from FHOPS 1.0.0
keep their labels (E1--E11); E6, E7, E8 (now E8a--E8c), E9 and E11 were
reformulated in 1.0.1, and E12, E13, INIT and OBJ2 are new. Section 2.3
keeps a prose summary of the objective and of each constraint family.
The main text is shorter, the appendix stays short, and the full
equations remain auditable in the version of the software the paper
describes.

\begin{center}\rule{0.5\linewidth}{0.5pt}\end{center}

\subsection*{Reviewer \#1}\label{reviewer-1}

\subsubsection{R1.1 ``All of my comments on the original version of the
article have been reviewed and corrected. Therefore, I recommend the
publication of this
paper.''}\label{r1.1-all-of-my-comments-on-the-original-version-of-the-article-have-been-reviewed-and-corrected.-therefore-i-recommend-the-publication-of-this-paper.}

\textbf{Response:} We thank Reviewer \#1 for the constructive
first-round comments and for recommending publication.

\begin{center}\rule{0.5\linewidth}{0.5pt}\end{center}

\subsection*{Reviewer \#2}\label{reviewer-2}

\subsubsection{R2.0 General comment: ``The authors have done an
excellent job. The application of software and process is well written
and explained. I would suggest adding market price of products based on
inventory which can give better insights on economic feasibility of the
harvesting
operations.''}\label{r2.0-general-comment-the-authors-have-done-an-excellent-job.-the-application-of-software-and-process-is-well-written-and-explained.-i-would-suggest-adding-market-price-of-products-based-on-inventory-which-can-give-better-insights-on-economic-feasibility-of-the-harvesting-operations.}

\textbf{Response:} We thank the reviewer for this positive assessment.
We address the market-price suggestion in detail under R2.9 below.

\subsubsection{R2.1 (title) ``The authors have done a good job
explaining FHOPS through the manuscript'' --- and R2.4 (contribution 1)
``Good
attempt''}\label{r2.1-title-the-authors-have-done-a-good-job-explaining-fhops-through-the-manuscript-and-r2.4-contribution-1-good-attempt}

\textbf{Response:} Thank you.

\subsubsection{R2.2 (keywords) ``Forest harvest operations and
optimisation already exist in the manuscript title consider another
wording.''}\label{r2.2-keywords-forest-harvest-operations-and-optimisation-already-exist-in-the-manuscript-title-consider-another-wording.}

\textbf{Response:} Agreed. We replaced the keywords that repeat the
title (``forest operations'', ``optimisation'') and replaced
``heuristics'' with more specific terms. The keywords are now:
\emph{machine scheduling; mixed-integer programming; metaheuristics;
reproducibility; forest management planning; sustainability}.

\subsubsection{R2.3 (abstract) ``Please expand''
{[}FHOPS{]}}\label{r2.3-abstract-please-expand-fhops}

\textbf{Response:} Done. The abstract now spells out the name at first
use: ``The Forest Harvesting Operations Planning System (FHOPS)
addresses that requirement\ldots{}''.

\subsubsection{R2.5 (Section 2) ``Does this include
trucking''}\label{r2.5-section-2-does-this-include-trucking}

\textbf{Response:} No.~The operational model schedules the in-block
machine roles up to and including loading at the landing. A loader works
only once a truckload (or the block's remaining volume, if smaller) is
staged at the landing (constraint block E9), but haul transport (truck
fleets, dispatching, and delivery to mills) is outside the model. We now
state this boundary explicitly in the opening paragraph of Section 2.

\subsubsection{R2.6 (Section 2) ``Helps in decision making to have a
coupled or decoupled harvesting
system''}\label{r2.6-section-2-helps-in-decision-making-to-have-a-coupled-or-decoupled-harvesting-system}

\textbf{Response:} We agree, and we now make this capability explicit in
Section 2.1. Each role can be given a head-start (in shifts of upstream
output), which sets the staged buffer B\_\{r,b\} in constraint block E8b
(the buffer is waived once the upstream roles have finished the block,
E8c):

\begin{itemize}
\tightlist
\item
  A positive head-start represents a decoupled system, e.g.~a roadside
  processor that starts only once a deck has accumulated.
\item
  A zero head-start approximates a coupled (hot) system, limited only by
  the one-shift staging lag of E7.
\end{itemize}

Users can therefore compare both configurations on the same scenario and
data contract. The reference ladder uses zero head-start for processing
and loading, apart from the loader's truckload threshold (E9).

\subsubsection{R2.7 (Section 2.1, production rates) ``I believe
utilization is considered here, because delays (operational/mechanical/
personal) has a major influence on the final cost of harvesting
operations''}\label{r2.7-section-2.1-production-rates-i-believe-utilization-is-considered-here-because-delays-operationalmechanical-personal-has-a-major-influence-on-the-final-cost-of-harvesting-operations}

\textbf{Response:} We agree that delays strongly affect production and
cost. FHOPS represents them at three levels, which we now describe in
the manuscript:

\begin{enumerate}
\def\labelenumi{\arabic{enumi}.}
\tightlist
\item
  \textbf{Production rates.} For the reference ladder, rates are derived
  from published productivity models and include each source model's
  delay or utilisation allowance (Section 2.1).
\item
  \textbf{Availability.} Machine availability flags in the shift
  calendar remove planned downtime, such as maintenance or crew days
  off, from the schedule (Section 2.1).
\item
  \textbf{Stochastic playback.} The solved schedule is replayed under
  sampled disruptions: machine downtime, weather, and landing shocks.
  The new text in Section 3.2 names these disturbance types and the
  sampling settings used for the reported results. Their effect on
  utilisation is shown in the playback figure.
\end{enumerate}

We also added a limitation to Section 4.4. Delay effects are currently
reported as utilisation and production losses. The schedule-level cost
outputs cover mobilisation only and do not yet convert delays into
standby or delay costs.

\subsubsection{R2.8 (Section 2.2) ``{[}re-solving{]} plays a key role
during weather uncertainities and
delays''}\label{r2.8-section-2.2-re-solving-plays-a-key-role-during-weather-uncertainities-and-delays}

\textbf{Response:} We agree. Section 2.2 now notes that FHOPS ships a
rolling-horizon driver (\texttt{fhops\ plan\ rolling}). It re-solves
successive planning windows with the SA or MILP solver while locking
near-term decisions, which supports periodic re-planning after weather
or delay disruptions. Before each window, the locked plan is replayed
and its state (remaining block volume, staged inventories, machine
positions, and locks) is carried into the next window; MILP windows add
an earliness tie-break so that they do not defer work beyond the locked
days. While preparing this answer we found that FHOPS 1.0.0 did not
carry this state between windows; the correction is part of FHOPS 1.0.1
(see the note at the end of this letter). The design of rolling-horizon
re-optimisation (planning-horizon length and re-optimisation frequency)
is evaluated in the companion study cited in the manuscript, across
three BC operating contexts and three problem sizes. Because the
rolling-horizon correction in FHOPS 1.0.1 affects that study's
experiments, the manuscript no longer quotes its quantitative findings;
it only points to the study.

\subsubsection{R2.9 (Section 2.3, mobilisation cost parameter) ``If
market values of various products can be added it helps to estimate the
revenue generated. This helps to get a better economic perspective of
the harvesting
operation.''}\label{r2.9-section-2.3-mobilisation-cost-parameter-if-market-values-of-various-products-can-be-added-it-helps-to-estimate-the-revenue-generated.-this-helps-to-get-a-better-economic-perspective-of-the-harvesting-operation.}

\textbf{Response:} We thank the reviewer for this suggestion and agree
that net revenue matters for the economic feasibility of harvesting
operations. However, it falls outside the decision scope of the model
presented here.

FHOPS addresses short-term (e.g.~multi-week), detailed operational
machine scheduling. At this level, the decisions that determine gross
revenue have already been made upstream, in tactical-operational
planning:

\begin{itemize}
\tightlist
\item
  which blocks are harvested;
\item
  with which harvest system and prescription;
\item
  and hence which assortments and volumes are produced.
\end{itemize}

The operational model takes these decisions as fixed inputs. It
optimises how the resulting workload is executed: machine--block
assignments, sequencing, timing, delivered production, mobilisation, and
landing use.

Because the assortment volumes are fixed upstream, product prices are
not needed to formulate or solve this scheduling problem. The objective
instead rewards delivering the planned volume and penalises volume left
unharvested. Economic feasibility in the revenue-and-margin sense is
evaluated at the tactical-operational level, where the
revenue-determining decisions are actually made.

We note that, since this manuscript was first submitted, FHOPS has
gained a very early-alpha tactical-operational planning layer, published
as pre-release FHOPS 1.1.0a1. It includes product values at mills and
terminals and profit or net-present-value objectives, which is the
planning level where the reviewer's suggestion applies. That layer is
provisional and outside the scope of this paper; Section 4.3 now
mentions it briefly.

To make this explicit to readers, we added a sentence to the opening
paragraph of Section 2. It positions the operational model below
tactical-operational planning and explains why the objective concerns
schedule execution rather than profit.

\subsubsection{R2.10 (Figure: deterministic vs.~stochastic utilisation)
``Please make the legends
readable''}\label{r2.10-figure-deterministic-vs.-stochastic-utilisation-please-make-the-legends-readable}

\textbf{Response:} Done. The overlapping title and clipped legend in the
original submission were already corrected in the first revision. In
this revision the figure is redrawn at the manuscript text width, so all
text, including the legend, prints at 9 pt or larger. The legend now
sits in a single row above the panels, the solver labels are in upper
case (SA, ILS), and the caption has been corrected. While regenerating
this figure we found and fixed defects in FHOPS's stochastic playback
events. The figure and the utilisation values in Section 3.2 were
regenerated on FHOPS 1.0.1 and have changed; see the note on FHOPS 1.0.1
at the end of this letter. The changes are small. Mean day-level
utilisation (deterministic → stochastic) is now 0.365 → 0.362 for tiny7
(previously 0.37 → 0.36), 0.566 → 0.562 for med42 with both SA and ILS
(previously 0.558 → 0.554 for SA and 0.564 → 0.559 for ILS), and 0.378 →
0.369 for synthetic-small (previously 0.62 → 0.60; its schedule changed
because the FHOPS 1.0.1 heuristics respect landing capacity). The
stochastic drop remains below 0.01 in every scenario and is still
largest for synthetic-small; SA and ILS now give identical deterministic
utilisation. Section 3.2 now reports all utilisation values to three
decimals so that these small differences are visible.

\begin{center}\rule{0.5\linewidth}{0.5pt}\end{center}

\subsection*{Note on FHOPS 1.0.1 (software
correction)}\label{note-on-fhops-1.0.1-software-correction}

The reviewers' questions on delays, re-solving, and figure legibility
(R2.7, R2.8, R2.10) led us to re-examine the corresponding FHOPS code
paths, and from there to audit FHOPS 1.0.0 more broadly: the MILP
formulation against the heuristics and playback, the heuristic
objective, and input validation. The audit found defects. We chose to
fix them and say so openly, rather than defer, work around, or silently
omit them. All fixes come with regression tests and are released as
\textbf{FHOPS 1.0.1}
(https://github.com/UBC-FRESH/fhops/releases/tag/v1.0.1;
\texttt{pip\ install\ fhops==1.0.1}), which the code metadata tables now
cite. The release notes list every change.

What was corrected:

\begin{enumerate}
\def\labelenumi{\arabic{enumi}.}
\tightlist
\item
  \textbf{Rolling-horizon driver.} State was not carried from one
  planning window to the next. Each window re-planned every block's full
  volume, including volume already delivered in earlier locked days;
  staged inventory between roles restarted at zero; and machine
  positions, user-specified locks, and blackout calendars were not
  carried across window boundaries. FHOPS 1.0.1 replays the locked plan
  before each window and carries its state forward. MILP windows also
  use an earliness tie-break, so they no longer defer work beyond the
  locked days.
\item
  \textbf{Stochastic playback.} Landing shocks were applied per
  assignment row rather than per calendar day, machine downtime always
  removed a whole shift instead of the configured duration, and weather
  and landing effects overwrote downtime instead of combining with it.
\item
  \textbf{MILP warm start.} Warm-starting the operational MILP failed
  with the default open-source solver (HiGHS). It now works, and FHOPS
  reports whether HiGHS accepted the start.
\item
  \textbf{Operational MILP formulation (Appendix A).} Several constraint
  blocks did not match the rules applied by the heuristics and by
  playback, so some MILP plans could not be replayed without sequencing
  violations or were charged differently:

  \begin{itemize}
  \tightlist
  \item
    staged inventories are now kept per upstream role (E7), so a role
    with several upstream roles processes only wood that each of them
    has staged;
  \item
    each role's output is capped at the volume the block still holds
    (E12), the head-start buffer is waived once the upstream roles have
    finished the block (E8c), and the loader threshold is the smaller of
    one truckload and the remaining volume (E9);
  \item
    landing capacity is enforced per shift and is hard at the default
    weight (E11); in 1.0.0 it was counted per day with a slack that cost
    nothing at the default weight, so it did not bind;
  \item
    machine moves are charged also when a machine idles between two
    blocks, and staying on a block is no longer charged (E6);
  \item
    the loader batching variables of 1.0.0, which imposed no
    restriction, were removed;
  \item
    timeline blackouts are enforced in the MILP (E1), and blocks without
    a harvest system carry no role obligations (E2).
  \end{itemize}
\item
  \textbf{Heuristic objective.} SA, ILS, and Tabu searched and reported
  a score that could be higher than the score of the schedule they
  returned, because the moves of machines not touched by the last repair
  were not charged. The reported objective is now a fresh evaluation of
  the returned schedule. The heuristics' repair step also respects
  landing capacity on every day, no longer starves downstream roles on
  capacity-limited landings, and penalises hard violations so that a
  schedule cannot gain from keeping an infeasible assignment.
\item
  \textbf{Validation and legacy code.} Scenario validation is stricter
  (e.g.~inconsistent locks, invalid initial states, and unknown
  harvest-system identifiers are rejected when a scenario is loaded).
  The legacy day-level MIP, which was infeasible for every scenario with
  a loader role, has been retired; \texttt{fhops\ solve-mip} now solves
  the operational MILP described in the paper.
\end{enumerate}

Effect on the manuscript:

\begin{itemize}
\tightlist
\item
  \textbf{All results were regenerated.} Every SoftwareX asset
  (benchmark, tuning, playback, costing, and scaling) was regenerated on
  FHOPS 1.0.1 in one pipeline run. Tables 4 and 5, the values quoted in
  Section 3, and the playback and scaling figures therefore changed. The
  main changes:

  \begin{itemize}
  \tightlist
  \item
    \emph{med42 and synthetic-small.} In the previous version their
    objectives were dominated by hard-violation penalties, mostly
    landing overloads that the FHOPS 1.0.0 heuristic repair step
    allowed. The 1.0.1 plans have no hard violations. The med42
    objectives are now positive and close (24130.32 for SA to 24329.17
    for ILS, previously −27271.22 to −46063.73), and ILS rather than SA
    has the best med42 score. Synthetic-small scores −39.79 for every
    solver (previously −75039.79), which is exactly its undelivered
    volume; it still delivers nothing and keeps its 39.79 m³ staged.
  \item
    \emph{Tuning (Table 5).} Δ against the SA default is 42.74 (tiny7),
    0.14 (small21), −990.50 (med42), and 0.00 (synthetic-small),
    previously 42.74, 10.23, −8214.08, and 0.00. Table 5 now has a
    Budget column and a note: each tuning run is much shorter than the
    SA default benchmark run it is compared with.
  \item
    \emph{Runtimes.} Every runtime in Table 4 is lower (by 16--39 \%),
    and the scaling runtimes are 20.69, 45.86, and 69.22 s for 4, 8, and
    16 blocks (previously 31.81, 84.88, and 108.97 s).
  \item
    \emph{Utilisation (Section 3.2).} See R2.10.
  \end{itemize}

  The main conclusions of Section 3 hold: objective, production,
  runtime, and mobilisation rank the solvers differently on med42, so
  solver choice should follow planning priorities; tuning gains are
  scenario-dependent and do not beat the SA default on med42;
  disturbance sensitivity is small and scenario-dependent; and runtime
  increases with the number of blocks. We revised the interpretations
  that no longer held. med42 is no longer described as
  penalty-dominated, and the synthetic-small score is now explained as
  unmet volume rather than as a penalty for infeasible schedules. The
  small21 tuning gain is now described as a tie, and Section 3.1
  explains that Δ compares a short tuning run with a much longer SA
  default run. The caption of Table 4 now says that only synthetic-small
  has a negative objective, and the sign-convention note in Section 2
  now names the unmet-volume term.
\item
  \textbf{Formulation.} The FHOPS 1.0.1 model, with labelled blocks and
  the list of changes from 1.0.0, is versioned with the software and
  summarised in Appendix A (see E.1). Section 2 describes the corrected
  behaviour (landing capacity per shift, rolling-horizon state, warm
  starts with HiGHS, exact heuristic objectives).
\item
  \textbf{Reproducibility.} Seeded heuristic results are
  bit-reproducible on a fixed platform, but last-bit floating-point
  differences between NumPy or Python builds can change an SA
  trajectory. Section 4.4 now says so, and the published assets record
  the platform used.
\item
  \textbf{Companion rolling-horizon study.} The companion study's
  rolling-horizon experiments were run with an earlier FHOPS version
  affected by the carry-forward defect. The manuscript therefore no
  longer states that study's quantitative conclusions (Sections 4.2 and
  5); it cites the study only as follow-on work built on FHOPS. Its
  experiments will be re-run on FHOPS 1.0.1.
\item
  \textbf{Synthetic tiers.} The FHOPS 1.0.1 synthetic generator produces
  scenarios whose machine roles do not match their harvest systems, so
  the synthetic tiers deliver no volume. Section 3 now states this
  explicitly; their rows are kept as runtime and reproducibility anchors
  only. The limitation is listed in the FHOPS 1.0.1 release notes and
  will be corrected in a later release.
\end{itemize}
